% Gesamttest «Lineare Funktionen» — Quelle des PDFs (python3 scripts/build-lp-pdf.py lineare-funktionen).
% Musterlösung und Raster: bewertungspaket.tex. Alle Werte mit python3 nachgerechnet (scripts/lp/lineare-funktionen/zahlen.py).
% Selbstcheck nach ~/TODO-BM-Vorkurs-LP03-15.md §2: zwei Wechsel Graph/Gleichung (G4 zeichnen, G5 ablesen) und
% Sachaufgaben mit der Infusion als Prüfungsanker (G6, G8). 25 Punkte, rund 30 Minuten. Nach der Prüfung vom 10.10.2026:
% neue Zahlen in G1, G3 bis G8 (keine Wiederholung von Kapitelaufgaben), G4 mit Nullstelle (K3), G2 in Prozentpunkten.
\documentclass[11pt]{article}
\usepackage{../lp-druck}
\usepackage{lf-gitter}
\renewcommand{\lpname}{Lineare Funktionen}
\renewcommand{\lpteilgebiet}{Funktionen ·}
\renewcommand{\lpfuss}{Gesamttest · 25 Punkte}
\begin{document}
\lpkopf{Gesamttest}{%
  \vspace{4pt}\noindent\begin{tabularx}{\linewidth}{@{}X X@{}}
  Code (kein Name): \hrulefill & Punkte: \hrulefill\ / 25
  \end{tabularx}}

\begin{lpinfo}{Anleitung}
\textbf{Zeit:} rund 30 Minuten. \textbf{Hilfsmittel:} Lineal, kein Taschenrechner.
Schreib zu jeder Aufgabe den \textbf{Rechenweg} auf; bewertet wird der Weg. Zeichne direkt in die Gitter.
Danach: selbst mit dem \emph{Bewertungspaket} bewerten oder die Lösung samt diesem Blatt scannen bzw.\ fotografieren
(ohne Namen und Standort) und mit dem Bewertungspaket einer KI geben.
\end{lpinfo}

\lpteil{Teil A · Anfangswert und Änderung}{Kapitel 1 · 6 Punkte}

\begin{lpaufgabe}{G1}{Werte}{3 P}
Gegeben ist $\;y = -3x + 12$.
\begin{enumerate}[label=(\alph*),leftmargin=*,itemsep=3pt]
\item Berechne $y$ für $x = -2$.
\item Für welches $x$ ist $y = -6$?
\end{enumerate}
\lpplatz{4}
\end{lpaufgabe}

\begin{lpaufgabe}{G2}{Akku}{3 P}
Ein Handyakku zeigt $80\,\%$. Pro Stunde Videoschauen sinkt die Anzeige um $6$ Prozentpunkte (auf $74\,\%$, dann $68\,\%$, \dots).
\begin{enumerate}[label=(\alph*),leftmargin=*,itemsep=3pt]
\item Schreib eine Gleichung für den Ladestand $y$ (in Prozent) nach $x$ Stunden.
\item Wie voll ist der Akku nach $5$ Stunden?
\item Nach wie vielen Stunden sind nur noch $20\,\%$ übrig?
\end{enumerate}
\lpplatz{5}
\end{lpaufgabe}

\lpteil{Teil B · Steigung, Graph und Gleichung}{Kapitel 2 bis 4 · 10 Punkte}

\begin{lpaufgabe}{G3}{Steigung}{3 P}
Berechne die Steigung der Geraden durch $A(-2 \mid 1)$ und $B(4 \mid 4)$. Steigt oder fällt die Gerade?
\lpplatz{3}
\end{lpaufgabe}

\begin{lpaufgabe}{G4}{Zeichnen}{4 P}
Zeichne die Gerade $\;y = -\tfrac{2}{3}x + 3$. Liegt $P(6 \mid -1)$ auf der Geraden? Wo schneidet die Gerade die $x$-Achse? Rechne beides nach.

\smallskip
\begin{minipage}[t]{0.55\linewidth}\vspace{0pt}
\gitter{-2}{9}{-4}{6}{}
\end{minipage}\hfill
\begin{minipage}[t]{0.42\linewidth}\vspace{0pt}
\lpplatz{6}
\end{minipage}
\end{lpaufgabe}

\begin{lpaufgabe}{G5}{Ablesen}{3 P}
Stelle die Gleichung der abgebildeten Geraden auf. Die beiden Punkte liegen auf Gitterpunkten.

\smallskip
\begin{minipage}[t]{0.45\linewidth}\vspace{0pt}
\gitter{-2}{6}{-4}{4}{\draw[lpblau,line width=1.2pt] (-2,2.5) -- (6,-3.5);
  \fill (0,1) circle (2.2pt); \fill (4,-2) circle (2.2pt);}
\end{minipage}\hfill
\begin{minipage}[t]{0.52\linewidth}\vspace{0pt}
\lpplatz{6}
\end{minipage}
\end{lpaufgabe}

\lpteil{Teil C · Gleichung aufstellen}{Kapitel 4 · 9 Punkte}

\begin{lpaufgabe}{G6}{Infusion}{3 P}
Ein Infusionsbeutel enthält $800$ ml. Die Infusion läuft mit $5$ ml pro Minute.
\begin{enumerate}[label=(\alph*),leftmargin=*,itemsep=3pt]
\item Schreib eine Gleichung für den Inhalt $V$ (in ml) nach $t$ Minuten.
\item Nach wie vielen Minuten ist der Beutel leer?
\item Zeichne den Graphen von $t = 0$ bis zum Zeitpunkt, an dem der Beutel leer ist.
\end{enumerate}

\smallskip
\begin{minipage}[t]{0.5\linewidth}\vspace{0pt}
\sachgitter{6}{8}{40}{100}{}
\end{minipage}\hfill
\begin{minipage}[t]{0.47\linewidth}\vspace{0pt}
\lpplatz{7}
\end{minipage}
\end{lpaufgabe}

\begin{lpaufgabe}{G7}{Zwei Punkte}{3 P}
Bestimme die Gleichung der Geraden durch $P(2 \mid -1)$ und $Q(6 \mid 5)$. Mach die Probe mit dem Punkt, den du nicht eingesetzt hast.
\lpplatz{5}
\end{lpaufgabe}

\begin{lpaufgabe}{G8}{Neue Rate}{3 P}
Ein Beutel enthält $700$ ml und läuft mit $6$ ml pro Minute. Nach $50$ Minuten wird auf $2$ ml pro Minute gestellt.
\begin{enumerate}[label=(\alph*),leftmargin=*,itemsep=3pt]
\item Wie viel ist beim Umstellen noch im Beutel?
\item Nach wie vielen Minuten (ab Beginn) ist er leer?
\item Zeichne den Graphen für den ganzen Ablauf.
\end{enumerate}

\smallskip
\begin{minipage}[t]{0.5\linewidth}\vspace{0pt}
\sachgitter{6}{7}{50}{100}{}
\end{minipage}\hfill
\begin{minipage}[t]{0.47\linewidth}\vspace{0pt}
\lpplatz{6}
\end{minipage}
\end{lpaufgabe}
\end{document}
